\begin{exercisechunk}
\item \textbf{A lock without observable progress.}
\label[exercise]{ex:l18-hidden-progress}
\exerciseprofile{3}{1}{2}{\exproof\quad\exinterpretation}{%
Identify the feedback assumption that makes stage-by-stage exploration
efficient.}
Fix integers $H\ge1$ and $n\ge0$. A hidden string $B$ is uniform on
$\{0,1\}^H$, and the learner has a random seed $U$ independent of $B$.
In query $i\in[n]$, the learner chooses a complete string
$X^i\in\{0,1\}^H$ and observes only
$Y_i=\I\{X^i=B\}$. The learner observes no intermediate states or
progress signals. Each $X^i$ may depend on all preceding strings and
responses and on $U$. After $n$ queries, the learner returns a string
$\widehat B$ determined by the query transcript and $U$.
\begin{enumerate}[(a)]
\item Prove that every such learner satisfies
\[
 \P(\widehat B=B)\le\min\left\{1,\frac{n+1}{2^H}\right\}.
\]
\item Deduce that $\P(\widehat B=B)\ge3/4$ requires
\[
 n\ge\left\lceil\frac34\,2^H\right\rceil-1.
\]
For comparison, suppose a query $X^i$ also returned the prefix indicators
$\I\{X^i_{1:k}=B_{1:k}\}$ for $k=1,\ldots,H$.
Explain how these indicators identify the first differing position in a
failed query and why the single terminal response $Y_i=0$ does not, in
general, supply that position.
\end{enumerate}
\begin{hints}
First consider a deterministic learner and the hypothetical transcript
in which every answer is zero.
\end{hints}
\end{exercisechunk}
